\documentclass[11pt]{article}
\usepackage{fontspec}
\usepackage[french,english]{babel}
\usepackage{amsmath}
\usepackage{amssymb}
\usepackage{amsthm}
\usepackage{mathrsfs}
\usepackage{tikz-cd}
\usepackage{enumitem}

\newcommand{\printedpage}[1]{\marginpar{\footnotesize p.~#1}}
\newcommand{\msp}[1]{\marginpar{\footnotesize [#1]}}

\begin{document}

\printedpage{113}
% head: A. Grothendieck : Dérivateurs, Ch. XII, Édition des Universités de Montpellier II et Paris VII
which do not factor through $F'$, the functor $T \longrightarrow T'$ being the functor of composition of $x \longrightarrow G$ with $G \xrightarrow{q} G'$
\[
\begin{tikzcd}[row sep=large]
F \arrow[r, hook] \arrow[d] & G \arrow[d] \\
F' \arrow[r, hook] & G'
\end{tikzcd}
\]
(NB The composite does not factor through $F'$, for otherwise, because of the cartesian character of $Q$, $x \longrightarrow G$ would factor through $F$.) But for given $x$, the datum of $x \longrightarrow G$ not factoring through $F$ is that of an element \msp{page 167} of $G(x)\backslash F(x)$, and because of the cocartesian character of $Q$,
\[
(G(x)\backslash F(x)) \longrightarrow (G'(x)\backslash F'(x))
\]
is bijective, i.e. ``it amounts to the same thing'' to give $x \longrightarrow G$ not factoring through $F$, or $x \longrightarrow G'$ not factoring through $F'$.

Thus we have proved the

\textbf{Theorem 5.} \textit{Let $X$ be in} Cat\textit{, and}
\[
Q : \qquad
\begin{tikzcd}[row sep=large, baseline=(current bounding box.center)]
F \arrow[r, hook, "i"] \arrow[d, "p"'] & G \arrow[d, "q"] \\
F' \arrow[r, hook, "i'"] & G'
\end{tikzcd}
\]
\textit{a cocartesian square in $X^\wedge$, with $i$} \texttt{[a]} \textit{monomorphism. Then $i'$ is} \texttt{[a]} \textit{monomorphism, the square $Q$ is also cartesian, and the corresponding square in} Cat$/X$ \textit{or in} Cat
\[
\begin{tikzcd}[row sep=large]
X/F \arrow[r, hook, "X/i"] \arrow[d, "X/p"'] & X/G \arrow[d, "X/q"] \\
X/F' \arrow[r, hook, "X/i'"] & X/G'
\end{tikzcd}
\]
\msp{page 168} \textit{is both cartesian and cocartesian. Moreover, if $W$ is a basic localizer in} Cat\textit{, satisfying} W(7 bis) (in addition to W(1,2,3))\textit{, then the square $X/Q$ in} Cat \textit{is $W$-cocartesian.}

\printedpage{114}
% head: A. Grothendieck : Dérivateurs, Ch. XII, Édition des Universités de Montpellier II et Paris VII
This last statement follows from what preceded and from corollary 2 of theorem 4, taking into account that $X/q$ is a local isomorphism, hence is $W$-smooth.

\textbf{Corollary:} (\footnote[111]{\textbf{NB} If we only want $i \in W_X \Longrightarrow i' \in W_X$, $p \in W_X \Longrightarrow q \in W_X$, it suffices to have W(7 bis) \textit{weak}, i.e. $i_0 \in W \Longrightarrow i \in W$ in the statement of W(7 bis), p. 76 (Loc(6$'$a) with the notation of XVI).}) \textit{We have} (assuming W(7 bis)) \textit{the equivalences}
\[
\begin{aligned}
i \in W_X \quad &\Longleftrightarrow \quad i' \in W_X \\
p \in W_X \quad &\Longleftrightarrow \quad q \in W_X,
\end{aligned}
\]
\textit{where}
\[
W_X \subseteq \mathit{Fl}(X^\wedge)
\]
\textit{is defined by (cf. Pursuing Stacks \ldots)}
\[
u \in W_X \qquad\qquad \Longleftrightarrow \qquad\qquad X/u : X/F \longrightarrow X/G \quad \textit{is in } W.
\]

Thus, in a category $X^\wedge$, amalgamated sums involving a monomorphism behave perfectly from the homotopy point of view, more precisely for the localizers $W_X$, provided \msp{page 169} only that $W$ satisfies W(7 bis). In particular, the monomorphisms in $X^\wedge$ are ``$W_X$-cofibrations'', in a sense similar to that of definition 2 c) (pages 99, 100 \texttt{[rather 109, 110]}).

\texttt{[page 170 does not exist.]}

\end{document}
